\begin{beginnerstatement}[In simple terms: FastAPI and HTTP]

The backend is the part that operates in the background; here, it is built using
the Python language and FastAPI. HTTP defines rules for how programs
transmit requests and responses over a network. An API is an agreed interface
through which the frontend communicates with the backend in this way.
A route or endpoint is an address within this API that can be called, with
\texttt{/api} grouping the shared check paths. As a liveness check,
\texttt{/api/health} reports whether the process is running; as a readiness
check, \texttt{/api/ready} reports whether it is ready for use and, when the
database is enabled, also checks its connection. An HTTP status code summarizes
the result as a number; here, HTTP~503 means that the service is temporarily not
ready. A container packages a service together with its runtime as a separate
unit, and the FastAPI backend runs in such a container during deployment. This
baseline does not yet include product-specific services or authentication.

\end{beginnerstatement}
